% Bagian: counterfactuals
% Bab: introduction
% Subbagian: paradoxes-material

\documentclass[../../../include/open-logic-section]{subfiles}

\begin{document}

\olfileid[id]{cnt}{int}{par}

\olsection{Paradoks Kondisional Material}

Salah seorang yang pertama mengkritik penggunaan $!A \lif !B$ sebagai cara
menyimbolkan pernyataan kondisional berbentuk ``jika \dots, maka \dots''
ialah C.~I. Lewis. Lewis mengkritik penggunaan kondisional material dalam
\emph{Principia Mathematica} karya Whitehead dan Russell, yang melafalkan
$\lif$ sebagai ``mengimplikasikan.'' Lewis dengan tepat mengeluhkan bahwa
jika $\lif$ berarti ``mengimplikasikan,'' maka setiap proposisi salah~$p$
mengimplikasikan bahwa $p$ mengimplikasikan~$q$, karena
$p \lif (p \lif q)$ benar jika~$p$ salah; dan setiap proposisi benar~$q$
mengimplikasikan bahwa $p$ mengimplikasikan~$q$, karena
$q \lif (p \lif q)$ benar jika~$q$ benar.

Para ahli logika tentu mengetahui bahwa \emph{implikasi}, yakni perikutan
logis, bukanlah konektif, melainkan relasi antara !!{formula} atau pernyataan.
Karena itu, untuk menghindari kebingungan, kita sebaiknya tidak membaca
$\lif$ sebagai ``mengimplikasikan.''\footnote{Pembacaan ``$\lif$'' sebagai
  ``mengimplikasikan'' masih tersebar luas di antara matematikawan dan
  ilmuwan komputer, meskipun para filsuf berusaha menghindari kebingungan
  yang ditunjukkan Lewis dengan melafalkannya sebagai ``hanya jika.''}
Selama kita tidak membacanya demikian, kekhawatiran khusus Lewis sama sekali
tidak muncul: $p$ tidak ``mengimplikasikan'' $q$ sekalipun kita menganggap
$p$ mewakili suatu kalimat bahasa alamiah yang salah. Untuk menentukan
apakah $p \Entails q$, kita harus mempertimbangkan \emph{semua}
!!{valuation}; dan $p \Entails/ q$ bahkan ketika kita menggunakan $p$ untuk
menyimbolkan suatu kalimat yang kebetulan salah.

Namun, tetap ada sesuatu yang ganjil tentang pernyataan ``jika \dots, maka
\dots'' seperti contoh Lewis berikut:
\begin{quote}
Jika bulan terbuat dari keju hijau, maka $2+2=4$.
\end{quote}
dan tentang inferensi-inferensi berikut:
\begin{quote}
  Bulan tidak terbuat dari keju hijau. Oleh karena itu, jika bulan terbuat
  dari keju hijau, maka $2+2=4$.

  $2+2 = 4$. Oleh karena itu, jika bulan terbuat dari keju hijau, maka
  $2+2=4$.
\end{quote}
Namun, jika ``jika \dots, maka \dots'' hanyalah $\lif$, kalimat tersebut
benar tanpa masalah dan inferensi-inferensi itu valid tanpa masalah.

Contoh lainnya menyangkut tautologi
$(!A \lif !B) \lor (!B \lif !A)$. Hal ini menyiratkan bahwa jika Anda
mengambil secara acak dua kalimat indikatif~$S$ dan~$T$ dari surat kabar,
kalimat ``Jika $S$, maka $T$, atau jika $T$, maka~$S$'' seharusnya benar.

\end{document}
